\section{Method}
\label{sec:method}

\subsection{Cross-domain task setup}
We adopt the bilingual Mongolian--Chinese cross-domain text detection
benchmark with two natural domains: handwritten archival pages
($H$, $1003$ images) and natural scene photographs ($S$, $1502$ images).
We study both transfer directions, which we denote $\hsource$
(handwritten as source, scene as target) and $\stoh$ (scene as source,
handwriting as target). The evaluation metric is polygon-level
$\hmean$ at IoU $\geq 0.5$.

\subsection{Detection backbone and training objective}
\label{sec:method:backbone}
We use DBNet~\cite{liao2020real} with a ResNet-18 backbone and an FPN
neck, taking $640 \times 640$ inputs. The network outputs a probability
map $p \in [0,1]^{H \times W}$ that scores per-pixel text presence. Ground
truth is constructed by shrinking each annotated polygon by a fixed offset
to a positive mask $g \in \{0,1\}^{H \times W}$. The detection loss is the
standard combination of dice and binary cross-entropy on $p$ versus $g$,
\begin{equation}
  \mathcal{L}_{\mathrm{det}}(p,g)
    = \mathcal{L}_{\mathrm{dice}}(p,g)
    + \lambda_{\mathrm{bce}}\,\mathcal{L}_{\mathrm{bce}}(p,g),
  \label{eq:det}
\end{equation}
with $\lambda_{\mathrm{bce}}=1.0$ throughout. At inference, $p$ is
binarised by a learned threshold and contours are extracted, then expanded
back by an unclip ratio of $1.5$. We use polygon-level $\hmean$ at
$\mathrm{IoU}\ge 0.5$ as the sole evaluation metric.

\subsection{Self-training UDA pipeline}
\label{sec:method:uda}
Our UDA pipeline combines (a) an exponential-moving-average (EMA)
teacher producing pseudo-labels on the unlabeled target domain, (b) an
image-level domain-adversarial classifier with a gradient-reversal layer
operating on fused FPN features (DANN~\cite{ganin2015dann}), and
(c) per-epoch validation on the target validation split, with
\texttt{best\_student.pth} saved at the highest target $\hmean$.
Let $\theta_s$ and $\theta_t$ denote the student and teacher parameters
respectively. The teacher is updated each step with an exponential moving
average,
\begin{equation}
  \theta_t \leftarrow \alpha\,\theta_t + (1-\alpha)\,\theta_s,
  \quad \alpha = 0.999.
  \label{eq:ema}
\end{equation}
For each minibatch, an image batch $x^{(s)}$ from the source domain
$\mathcal{D}_s$ (with ground-truth masks $g^{(s)}$) and an unlabeled
batch $x^{(t)}$ from the target domain $\mathcal{D}_t$ are drawn.
Pseudo-labels $\tilde g^{(t)}$ on the target batch are produced by
binarising the teacher's probability map $p_t(x^{(t)})$ at a confidence
threshold $\tau_{\mathrm{conf}}=0.55$. A per-image keep-flag
\begin{equation}
  m_i =
  \mathbf{1}\!\left[\,
    \max_{u,v} p_t\bigl(x_i^{(t)}\bigr)_{uv} > \tau_{\mathrm{conf}}
  \right]
  \label{eq:keep}
\end{equation}
discards images on which the teacher is too uncertain to provide any
pseudo-signal; the realised fraction $\mathbb{E}[m_i]$ is the
\emph{pseudo coverage} reported in our logs. The total student
objective at training step $k$ is
\begin{equation}
\begin{aligned}
  \mathcal{L}_{\mathrm{stud}}
   = \;& \mathcal{L}_{\mathrm{det}}\!\bigl(p_s(x^{(s)}), g^{(s)}\bigr) \\
     & + w_p(k)\;\overline{m_i\,
        \mathcal{L}_{\mathrm{det}}\!\bigl(p_s(x_i^{(t)}),\,\tilde g_i^{(t)}\bigr)} \\
     & + w_d(k)\;\mathcal{L}_{\mathrm{dom}}\!\bigl(F(x^{(s)}),F(x^{(t)})\bigr),
\end{aligned}
  \label{eq:total}
\end{equation}
where $F(\cdot)$ denotes the FPN-fused feature map fed into the domain
classifier, $\mathcal{L}_{\mathrm{dom}}$ is binary cross-entropy on the
source/target label, and $w_p, w_d$ are the pseudo-label and
domain-loss weights described in Sec.~\ref{sec:method:weighting}. The
overline denotes the per-batch mean.

The GRL coefficient on the domain branch follows the canonical schedule
$\gamma(k) = 2/(1+\exp(-10\,k/K)) - 1$, where $K$ is the total number
of training steps; $\gamma$ saturates to $1.0$ in the second half of
training, making the domain head progressively harder to fool.

\subsection{Cold-start initialisation (key contribution)}
\label{sec:method:coldstart}
A naive initialisation -- random weights for the projection layers plus
ImageNet-pretrained backbone -- is insufficient when the source-only
model's accuracy on the target domain is near zero. The teacher,
inheriting the same weights, generates pseudo-labels that are
indistinguishable from random prediction, and the student converges to
the trivial fixed point of predicting nothing on target images. We
observe this failure first hand on the $\stoh$ direction: a 40-epoch
training run from random initialisation collapses to $\hmean = 0.0029$,
with the domain head loss saturating to zero (no useful gradient signal)
and pseudo-label coverage stuck at $1.0$ on garbage labels.

We propose a minimal remedy. Both the student and the teacher are
initialised from the source-only checkpoint
(\texttt{best.pth} on the source domain), so that the first round of
pseudo-labels is generated by a model that has learned generic text
features (even if its target-domain precision is still very low). This
gives the EMA teacher a non-trivial trajectory to follow as the student
improves.

\subsection{Conservative pseudo-label weighting}
\label{sec:method:weighting}
The pseudo-label and domain weights both ramp linearly from $0$ at
the start of training. The pseudo-weight rises over the first
$w_p^{\mathrm{warm}}$ epochs to $w_p^{\max}$ and is held constant
afterwards; the domain weight follows the same ramp with maximum
$w_d^{\max}=0.1$. Formally,
\begin{equation}
  w_p(k) \;=\; w_p^{\max}\,\min\!\Bigl(1,\;\frac{k}{w_p^{\mathrm{warm}}\,S}\Bigr),
  \label{eq:ramp}
\end{equation}
where $S$ is the number of steps per epoch.

The cap $w_p^{\max}$ is the single most consequential hyperparameter
in our pipeline. Setting $w_p^{\max}=0.5$ (our v1 ablation) reproduces
the typical literature value and improves early epochs, but its
sensitivity to teacher drift is so high that the model collapses by
epoch 5. Setting $w_p^{\max}=0.1$ (v3) preserves the early-epoch
gain while delaying the collapse onset to epoch~9 in $\hsource$ and
epoch~2 in $\stoh$. The slower onset is what allows the
\texttt{best\_student} checkpoint policy to capture useful weights at
all. Sec.~\ref{sec:exp:long} quantifies the collapse and motivates
the conservative cap as a structural choice rather than a budget
compromise.
